% !TeX root = ../main.tex
\section{Discussion}
\label{sec:discussion}

The main implication is not that one analog precision is universally optimal. It is that analog deployment has to be staged: first prevent the network from overfitting a single hardware instance, then choose readout and memory precision under the target device physics.

The first deployment frontier is algorithmic. Standard fixed-mask hardware-aware training can be effective when precision is sufficient or when the trained hardware instance is reused, but it becomes fragile at extreme low precision because the optimizer can exploit one sampled mismatch field. The model then performs well for the training instance and poorly for fresh hardware draws. Ensemble HAT changes the objective: by exposing the model to a distribution of mismatch fields during training, it discourages instance-specific solutions and makes fresh-instance transfer part of the optimization target.

That algorithmic rescue does not remove the device frontier; it makes the frontier measurable. Once cross-instance overfitting is controlled, the limiting question becomes whether the memory substrate maintains a usable state after programming. In the tested PCM UnitCell models, 8-bit precision is the strongest practical operating point because it combines high fresh accuracy (77.60\%) with drift stability. The 6-bit result is more informative than a monotonic scaling law would be: it is drift-stable, but D2D perturbations reduce mean accuracy (68.44\%) and increase seed-to-seed variance (std 6.03\%). It therefore marks a transition zone between the quantization-dominated 4-bit regime and the noise-absorbed 8-bit regime.

The joint ADC--D2D sweep gives a similarly ordered hardware roadmap. ADC resolution is a hard gate: below 6 bits, classification accuracy collapses regardless of device quality. Above that threshold, device-to-device variability becomes the dominant accuracy limiter. This ordering matters for hardware co-design. Improving device uniformity before the readout interface clears the ADC cliff will not recover the model; after the cliff is cleared, mismatch reduction becomes the correct target.

The conclusions are bounded by the present behavioral simulation stack. Device profiles are calibrated from literature values or a single fabrication batch, so the numbers should not be read as cross-batch, temperature, aging, or circuit-layout coverage. We model quantization, D2D/C2C variability, ADC precision, write nonlinearity, and retention, but defer parasitic IR drop, sneak paths, and pulse-faithful programming to follow-on hardware-calibrated studies.

The energy estimates are likewise first-order analytical projections from literature-proxy constants ($E_{\text{cell}}=100$~fJ, $E_{\text{ADC},8b}=25$~fJ, $E_{\text{DAC},8b}=30$~fJ; Methods). They rank deployment risks rather than predict silicon performance. The current ablation matrix also prioritizes non-ideality coverage over dataset breadth; cross-architecture TinyImageNet sweeps (see Supplementary Information) are treated as validation rather than as main-text claims.

\begin{table}[htbp]
\centering
\caption{\textbf{Deployment Risk Ranking and Mitigation Strategies.} Factors are ranked by their impact on categorical accuracy as identified by Sobol sensitivity analysis and precision-ladder sweeps.}
\label{tab:deployment_risks}
\begin{tabular}{lll}
\toprule
\textbf{Factor} & \textbf{Impact Rank} & \textbf{Mitigation Strategy} \\
\midrule
ADC Resolution & High (Hard Gate) & Secure $\geq$6-bit readout precision \\
Stochastic D2D & High (Generalization) & Ensemble HAT (distribution matching) \\
PCM Drift & Moderate (Time-bound) & 8-bit for stability; refresh 4-bit periodically \\
Write Nonlinearity & Moderate (Optimization) & STE gradient scaling \\
Shot Noise & Low (Transduction) & Front-end inverse-gamma compensation \\
\bottomrule
\end{tabular}
\end{table}

The next tests are straightforward: replace UnitCell presets with measured-array statistics, verify whether the non-monotonic precision-accuracy curve and 6-bit transition-zone behavior persist under measured devices, and evaluate memory-bound inference components such as analog KV-cache with explicit train/eval noise separation. These extensions do not alter the present conclusion: algorithmic cross-instance resampling and physical precision-retention selection are separate deployment decisions.
